\chapterimage{figs/chapterhead.png}
\chapter{Parser Architecture}
\label{chap:parser-architecture}

This chapter documents the current parser subsystem as it exists in the
repository today. The parser is not treated as a generic text utility; it is a
model-aware service that translates expressions, function calls, assignments,
and selected control forms into numeric values and model mutations.

\noindent\textbf{Objective.} Explain how the parser is wired, how the scanner
and grammar cooperate, how model context influences evaluation, and which
parts of the regeneration workflow are still incomplete.

\noindent\textbf{Audience.} Developers who need to change parsing,
expression evaluation, model-language semantics, or parser-related build
machinery.

\noindent\textbf{Scope.} Cover the current \texttt{Parser\_if} contract, the
\texttt{genesyspp\_driver} implementation, the Bison/Flex source files, the
generated parser artifacts, model-backed symbol resolution, error handling,
ownership, regeneration, and the available tests.

\noindent\textbf{Status.} Confirmed by code inspection and unit tests. The
current parser grammar, scanner, and driver are active; optional regeneration
is present in CMake; the higher-level \texttt{ParserManager} generation flow is
still a stub and must not be read as a finished feature.

The chapter follows Figure~\ref{fig:parser-pipeline},
Figure~\ref{fig:parser-model-dependency},
Figure~\ref{fig:parser-generation-flow}, and
Figure~\ref{fig:parser-error-flow}.

\section{What the parser service is}
\label{sec:parser-service}

The public contract is intentionally small. \texttt{Parser\_if} exposes two
expression-evaluation forms, an error-message accessor, sampler injection, and
access to the internal driver for advanced use cases. That is the architectural
signal that matters most: parser access is a simulator service, not an isolated
library API.

The interface currently supports:
\begin{itemize}
\item expression evaluation that may throw;
\item expression evaluation that reports success and an error string;
\item access to the last parser error message;
\item association with a \texttt{Sampler\_if} instance for stochastic terms;
\item exposure of the internal \texttt{genesyspp\_driver}.
\end{itemize}

This contract is used by the kernel and by model-facing code. The current
implementation is centered on \texttt{genesyspp\_driver}, which owns the parse
session state, the selected model, the sampler pointer, error text, result
value, and the registry used to track referenced data elements.

\section{Source text to value}
\label{sec:parser-source-to-value}

The concrete parser lives under \texttt{source/parser/}. The important files
are:
\begin{itemize}
\item \texttt{Genesys++-driver.h} and \texttt{Genesys++-driver.cpp} for the
      parse-session wrapper;
\item \texttt{parserBisonFlex/bisonparser.yy} for the grammar;
\item \texttt{parserBisonFlex/lexerparser.ll} for the scanner;
\item \texttt{GenesysParser.h}, \texttt{GenesysParser.cpp}, and
      \texttt{Genesys++-scanner.cpp} for the committed generated artifacts;
\item \texttt{CMakeLists.txt} for the parser target and optional regeneration.
\end{itemize}

The path from source text to numeric value has three stages:
\begin{enumerate}
\item the scanner converts characters into typed tokens;
\item the Bison grammar reduces those tokens into expressions, commands, and
      model-aware function calls;
\item the driver stores the resulting value, error state, and side effects.
\end{enumerate}

The scanner recognizes numbers, operators, booleans, math functions,
probabilistic functions, simulation symbols, kernel helpers, and plugin
extensions. Identifiers are not treated as plain text; the scanner resolves
them against the current model first. It checks attributes, statistics
collectors, counters, simulation responses, simulation controls, and then the
known plugin-backed data definitions.

That resolution step is why the parser is model-aware. The scanner already
depends on the live model state before the grammar even starts reducing.

\begin{figure}[htbp]
\centering
% FIGURE-SPEC-BEGIN
% ID: fig:parser-pipeline
% Source of truth:
%   - source/parser/parserBisonFlex/lexerparser.ll
%   - source/parser/parserBisonFlex/bisonparser.yy
%   - source/parser/Genesys++-driver.h
%   - source/parser/Genesys++-driver.cpp
% Purpose:
%   Show the current path from source text through tokens and grammar to the final value or error.
% Required visual content:
%   - source text input
%   - scanner/token layer
%   - grammar reductions
%   - driver state and returned value
%   - explicit error branch
% Accessibility description:
%   A pipeline diagram that lets the reader follow one expression from characters to evaluated result.
% Validation:
%   The diagram must match the current scanner, grammar, and driver responsibilities.
% FIGURE-SPEC-END
\figureplaceholder{figs/generated/parser-pipeline}
\caption{Current parser pipeline from source text to value.}
\label{fig:parser-pipeline}
\end{figure}

Figure~\ref{fig:parser-pipeline} should be read as the operational view of
the subsystem. The scanner is not a passive preprocessor, and the grammar is
not a generic calculator grammar. Both are coupled to model-aware semantics.

\section{Driver and model dependency}
\label{sec:parser-driver-model}

The driver is the bridge between the grammar and the rest of the kernel. It
stores the current model pointer, sampler pointer, parse string or file name,
error state, the final numeric result, and an optional registry of referenced
data elements. The copy and move operations are implemented explicitly, with a
deep copy of the referenced-data registry so copied drivers do not share list
ownership.

That ownership detail matters. The registry is heap-allocated, and the driver
cleans it up in the destructor. The copy and move paths therefore preserve the
driver as a self-contained parse-session object rather than a shallow handle.

The model dependency is also explicit in the helper functions exposed by the
driver:
\begin{itemize}
\item \texttt{findSimulationResponse()};
\item \texttt{findSimulationControl()};
\item \texttt{getSimulationResponseValueAsDouble()};
\item \texttt{getSimulationControlValueAsDouble()};
\item \texttt{traceWarning()}.
\end{itemize}

These helpers are used when the grammar evaluates model symbols that are not
plain arithmetic values. A simulation response or control is looked up by
name, converted to \texttt{double} when possible, and traced as a warning if
its stored value is non-numeric.

\begin{figure}[htbp]
\centering
% FIGURE-SPEC-BEGIN
% ID: fig:parser-model-dependency
% Source of truth:
%   - source/kernel/simulator/Parser_if.h
%   - source/parser/Genesys++-driver.h
%   - source/parser/Genesys++-driver.cpp
%   - source/parser/parserBisonFlex/lexerparser.ll
%   - source/parser/parserBisonFlex/bisonparser.yy
% Purpose:
%   Show how the parser depends on the live model, sampler, and kernel services.
% Required visual content:
%   - Parser_if contract
%   - genesyspp_driver as session object
%   - Model and Simulation access
%   - Sampler_if dependency for stochastic functions
%   - helper lookups for responses, controls, and data definitions
% Accessibility description:
%   A dependency diagram that makes the runtime coupling to the model explicit.
% Validation:
%   The diagram must reflect the actual helper methods and data lookups in the current code.
% FIGURE-SPEC-END
\figureplaceholder{figs/generated/parser-model-dependency}
\caption{Current model and kernel dependencies used by the parser.}
\label{fig:parser-model-dependency}
\end{figure}

Figure~\ref{fig:parser-model-dependency} makes the central point visible:
the parser evaluates expressions in model context, not in isolation. This is
why parse errors and value lookups are tied to current simulation state.

\section{Grammar and semantics}
\label{sec:parser-grammar-semantics}

The grammar in \texttt{bisonparser.yy} is a typed Bison C++ grammar using the
\texttt{lalr1.cc} skeleton. It defines the current expression language through
precedence levels for logical, relational, additive, multiplicative, power,
unary, and primary forms. It also defines assignment, command forms, kernel
helpers, math functions, probability functions, plugin-backed functions, and
model-backed symbols.

The current grammar covers:
\begin{itemize}
\item arithmetic, relational, and logical operators;
\item conditional forms with \texttt{if};
\item basic loop syntax with \texttt{for};
\item kernel functions such as \texttt{TNOW}, \texttt{TFIN}, \texttt{MAXREP},
      \texttt{NUMREP}, \texttt{IDENT}, and \texttt{entitiesWIP};
\item math and trigonometric functions;
\item stochastic functions such as \texttt{expo}, \texttt{norm},
      \texttt{unif}, \texttt{weib}, \texttt{logn}, \texttt{gamm},
      \texttt{erla}, \texttt{tria}, \texttt{beta}, and \texttt{rnd};
\item model symbols for attributes, simulation responses, simulation controls,
      variables, formulas, queues, resources, sets, and counters.
\end{itemize}

The most important semantic decision is that grammar actions can mutate the
model. Assignments to attributes and variables do not just return a numeric
value; they also write back to the current entity or to the data-definition
object. The parser therefore acts as both evaluator and model mutator.

Some details remain intentionally conservative:
\begin{itemize}
\item \texttt{FORM} still returns a placeholder numeric value instead of
      re-entering the parser recursively;
\item \texttt{fDISC}, \texttt{fLASTINQ}, and \texttt{fRESSEIZES} are present
      in the grammar but are not fully implemented;
\item a few current actions still rely on the current event and current entity
      being present, so those paths must be handled with care.
\end{itemize}

Those gaps are architectural debt, not hidden success conditions. The chapter
records them so the reader does not infer support that the current code does
not actually provide.

\begin{figure}[htbp]
\centering
% FIGURE-SPEC-BEGIN
% ID: fig:parser-generation-flow
% Source of truth:
%   - source/parser/CMakeLists.txt
%   - source/parser/parserBisonFlex/bisonparser.yy
%   - source/parser/parserBisonFlex/lexerparser.ll
%   - source/parser/GenesysParser.cpp
%   - source/parser/GenesysParser.h
%   - source/parser/Genesys++-scanner.cpp
% Purpose:
%   Show how the committed generated files and the optional regeneration path are related.
% Required visual content:
%   - Bison/Flex source files
%   - optional regeneration toggle in CMake
%   - generated parser headers and sources
%   - static library assembly
%   - note that ParserManager generation is still a stub
% Accessibility description:
%   A build-flow diagram that distinguishes committed generated code from optional regeneration.
% Validation:
%   The diagram must match the current CMake option, generated artifacts, and current stubbed generation API.
% FIGURE-SPEC-END
\figureplaceholder{figs/generated/parser-generation-flow}
\caption{Current parser generation and rebuild flow.}
\label{fig:parser-generation-flow}
\end{figure}

Figure~\ref{fig:parser-generation-flow} is important because it separates two
things that are easy to confuse. The repository already ships generated parser
artifacts, but the regeneration workflow is only optionally enabled by CMake
and the higher-level \texttt{ParserManager} API is still unfinished.

\section{Error handling and limits}
\label{sec:parser-errors-limits}

Error handling is split between the lexer, the grammar, and the driver.
\texttt{ILLEGAL} tokens are turned into explicit parser failures. The driver
also catches thrown \texttt{std::exception}, thrown \texttt{std::string}, and
unknown exceptions in \texttt{parse\_file()} and \texttt{parse\_str()}. When
that happens, it stores an error message, sets the result to \texttt{-1}, and
either rethrows or returns depending on the configured exception mode.

The grammar also distinguishes between deterministic evaluation and warning
paths. For example, non-numeric simulation-response or simulation-control
values are converted conservatively to \texttt{0.0} after a trace warning. That
behavior is intentionally pragmatic, but it should not be confused with a full
validation guarantee.

\begin{figure}[htbp]
\centering
% FIGURE-SPEC-BEGIN
% ID: fig:parser-error-flow
% Source of truth:
%   - source/parser/parserBisonFlex/lexerparser.ll
%   - source/parser/parserBisonFlex/bisonparser.yy
%   - source/parser/Genesys++-driver.cpp
% Purpose:
%   Show how illegal tokens, grammar failures, and driver exceptions become reported parser errors.
% Required visual content:
%   - illegal token branch
%   - parse_str/parse_file exception handling
%   - error message storage
%   - result reset to -1
%   - warning path for non-numeric simulation symbols
% Accessibility description:
%   A flow diagram that makes the current error-reporting behavior easy to distinguish from successful evaluation.
% Validation:
%   The diagram must match the actual catch blocks and error-message handling in the driver.
% FIGURE-SPEC-END
\figureplaceholder{figs/generated/parser-error-flow}
\caption{Current error flow for lexing, parsing, and evaluation.}
\label{fig:parser-error-flow}
\end{figure}

Figure~\ref{fig:parser-error-flow} should make the current failure behavior
obvious. A parse failure is not just a syntax error in the abstract; it is a
specific path through the lexer, grammar, and driver that ends with a stored
message and a negative result.

\section{Regeneration and build integration}
\label{sec:parser-regeneration}

The parser build is controlled by \texttt{source/parser/CMakeLists.txt}. By
default, the static \texttt{genesys\_parser} library is built from the checked
in \texttt{*.cpp} files in \texttt{source/parser/}. An optional
\texttt{GENESYS\_PARSER\_REGENERATE} switch enables Bison and Flex, which then
regenerate the parser and scanner sources from
\texttt{parserBisonFlex/bisonparser.yy} and
\texttt{parserBisonFlex/lexerparser.ll}.

This is a useful but limited workflow:
\begin{itemize}
\item the regeneration path is present and explicit;
\item the project does not require regeneration for every build;
\item the current generation API in \texttt{ParserManager} is only a declared
      extension point;
\item the chapter should not imply a fully automated parser-creation pipeline.
\end{itemize}

The practical consequence is that parser changes usually follow two distinct
steps:
\begin{enumerate}
\item update the grammar/scanner and, if needed, regenerate the committed
      artifacts;
\item rebuild the parser library and run the expression tests.
\end{enumerate}

That sequence keeps the checked-in generated files and the source grammar in
sync without pretending that the project already owns a dynamic regeneration
service.

\section{Tests and safe modification}
\label{sec:parser-tests}

The current test surface confirms the parser in a focused way. The most useful
checks are:
\begin{itemize}
\item \texttt{test\_parser\_expressions.cpp}, which exercises precedence,
      right-associative power, logical operators, math functions, indexed
      variable access, indexed attribute access, assignment, and event-bound
      evaluation;
\item \texttt{test\_support\_parsermanager.cpp}, which confirms the default
      construction state of \texttt{ParserManager} result and path holders.
\end{itemize}

Those tests are enough to anchor the current chapter text, but they do not
prove broad language completeness. In particular, they do not resolve the
remaining stubs around formula evaluation, discrete sampling, or parser
regeneration.

When modifying this subsystem, the safe sequence is:
\begin{enumerate}
\item inspect \texttt{bisonparser.yy}, \texttt{lexerparser.ll}, and the driver
      together;
\item verify whether the affected symbol is a scanner token, a grammar
      production, or a driver helper;
\item check whether the change requires regeneration or only a source edit;
\item run the parser-focused tests before widening validation.
\end{enumerate}

That workflow keeps expression semantics, model coupling, and error behavior
aligned. It also prevents the chapter from overstating the maturity of the
unfinished regeneration layer.
